\documentclass[aspectratio=169,8pt]{beamer}
\usetheme{Madrid}
\usepackage[utf8]{inputenc}
\usepackage[T1]{fontenc}
\usepackage{amsmath,amssymb}
\usepackage{booktabs}
\usepackage{enumitem}
\setlist[enumerate]{label=\Alph*.,leftmargin=*,itemsep=1pt,topsep=2pt}
\setlist[itemize]{leftmargin=*,itemsep=1pt,topsep=2pt}
\setbeamerfont{block body}{size=\tiny}
\setbeamerfont{block title}{size=\scriptsize}
\setbeamerfont{frametitle}{size=\footnotesize}
\setbeamerfont{normal text}{size=\tiny}
\setbeamertemplate{navigation symbols}{}
\setbeamertemplate{itemize item}{\tiny$\bullet$}
\setbeamertemplate{itemize subitem}{\tiny$\circ$}

\begin{document}
	
	\begin{frame}{Q550 --- ROC Curve}
		\begin{block}{Question}
			An ROC curve plots:
			\begin{enumerate}
				\item Precision against Recall
				\item True Positive Rate against False Positive Rate
				\item Accuracy against Threshold
				\item Recall against Precision
			\end{enumerate}
		\end{block}
		
		\begin{exampleblock}{Answer: B}
			True Positive Rate against False Positive Rate
		\end{exampleblock}
		
		\begin{block}{Explanation}
			An \textbf{ROC curve} (Receiver Operating Characteristic) plots the
			\textbf{True Positive Rate (TPR)} against the \textbf{False Positive Rate (FPR)}
			across all possible classification thresholds:
			\[
			\text{TPR} = \frac{TP}{TP+FN} = \text{Sensitivity / Recall}
			\]
			\[
			\text{FPR} = \frac{FP}{FP+TN} = 1 - \text{Specificity}
			\]
			
			\textbf{How to read it:}
			\begin{itemize}
				\item The top-left corner (TPR = 1, FPR = 0) is perfect classification.
				\item The diagonal line (TPR = FPR) represents random guessing.
				\item A curve closer to the top-left corner is a better model.
			\end{itemize}
			
			\textbf{Why the others are wrong:}
			\begin{itemize}
				\item \textbf{A:} Precision vs. Recall is the \emph{Precision-Recall curve},
				a different plot (often better for imbalanced data).
				\item \textbf{C:} Accuracy vs. Threshold is a threshold-sweep plot, not ROC.
				\item \textbf{D:} Recall vs. Precision is not a standard ROC curve.
			\end{itemize}
		\end{block}
	\end{frame}
	
	\begin{frame}{Q551 --- AUC Interpretation}
		\begin{block}{Question}
			What does AUC primarily measure?
			\begin{enumerate}
				\item Overall prediction accuracy
				\item Calibration quality
				\item Ability to distinguish positives from negatives
				\item Number of correct classifications
			\end{enumerate}
		\end{block}
		
		\begin{exampleblock}{Answer: C}
			Ability to distinguish positives from negatives
		\end{exampleblock}
		
		\begin{block}{Explanation}
			\textbf{AUC} (Area Under the ROC Curve) measures \textbf{discrimination}---the
			model's ability to \emph{separate} positives from negatives.
			
			\textbf{Probabilistic interpretation:} AUC is the probability that a randomly
			chosen positive example receives a higher predicted score than a randomly
			chosen negative example.
			
			\textbf{Why the others are wrong:}
			\begin{itemize}
				\item \textbf{A:} AUC is not accuracy. A model can have high AUC but low
				accuracy at a specific threshold (and vice versa).
				\item \textbf{B:} Calibration measures whether predicted probabilities match
				observed frequencies---a different concept from discrimination.
				\item \textbf{D:} AUC is threshold-independent and does not count correct
				classifications directly.
			\end{itemize}
			
			\textbf{Key point:} AUC is about \emph{ranking}, not about correctness at a
			single operating point.
		\end{block}
	\end{frame}
	
	\begin{frame}{Q552 --- Random Classifier}
		\begin{block}{Question}
			An AUC of 0.50 suggests:
			\begin{enumerate}
				\item Perfect classification
				\item Random guessing
				\item Severe overfitting
				\item High precision
			\end{enumerate}
		\end{block}
		
		\begin{exampleblock}{Answer: B}
			Random guessing
		\end{exampleblock}
		
		\begin{block}{Explanation}
			An AUC of \textbf{0.50} means the ROC curve lies on the diagonal line
			(TPR = FPR). The model has \textbf{no discrimination ability}---it performs
			exactly like random guessing.
			
			\textbf{AUC scale:}
			\begin{itemize}
				\item 0.50 = random guessing (no discrimination)
				\item 0.50--0.70 = poor
				\item 0.70--0.80 = acceptable
				\item 0.80--0.90 = excellent
				\item 0.90--1.00 = outstanding
				\item 1.00 = perfect separation
			\end{itemize}
			
			\textbf{Why the others are wrong:}
			\begin{itemize}
				\item \textbf{A:} Perfect classification corresponds to AUC = 1.00.
				\item \textbf{C:} AUC = 0.50 says nothing about overfitting; it means no
				signal.
				\item \textbf{D:} AUC does not measure precision.
			\end{itemize}
		\end{block}
	\end{frame}
	
	\begin{frame}{Q553 --- Perfect Classifier}
		\begin{block}{Question}
			Which AUC value corresponds to a perfect classifier?
			\begin{enumerate}
				\item 0.50
				\item 0.75
				\item 0.90
				\item 1.00
			\end{enumerate}
		\end{block}
		
		\begin{exampleblock}{Answer: D}
			1.00
		\end{exampleblock}
		
		\begin{block}{Explanation}
			An AUC of \textbf{1.00} means the ROC curve passes through the top-left corner
			(TPR = 1, FPR = 0). The model perfectly separates positives from negatives at
			some threshold.
			
			\textbf{Why the others are wrong:}
			\begin{itemize}
				\item \textbf{A (0.50):} Random guessing.
				\item \textbf{B (0.75):} Acceptable but not perfect.
				\item \textbf{C (0.90):} Excellent but not perfect.
			\end{itemize}
			
			\textbf{Note:} In practice, AUC = 1.00 usually indicates a data leak, label
			leakage, or an overly easy problem---real-world models rarely achieve it.
		\end{block}
	\end{frame}
	
	\begin{frame}{Q554 --- Threshold Independence}
		\begin{block}{Question}
			Why is AUC often preferred to accuracy?
			\begin{enumerate}
				\item It avoids the need for labels
				\item It evaluates performance across all thresholds
				\item It measures clustering quality
				\item It measures calibration
			\end{enumerate}
		\end{block}
		
		\begin{exampleblock}{Answer: B}
			It evaluates performance across all thresholds
		\end{exampleblock}
		
		\begin{block}{Explanation}
			Accuracy depends on a \textbf{single threshold} (usually 0.5). AUC summarizes
			performance across \textbf{all possible thresholds}, so it is
			\textbf{threshold-independent}.
			
			\textbf{Why this matters:}
			\begin{itemize}
				\item You can compare models without committing to a threshold.
				\item AUC is robust to class imbalance (unlike accuracy).
				\item It reflects the model's ranking quality overall.
			\end{itemize}
			
			\textbf{Why the others are wrong:}
			\begin{itemize}
				\item \textbf{A:} AUC still requires ground-truth labels.
				\item \textbf{C:} Clustering is unsupervised; AUC is for classification.
				\item \textbf{D:} Calibration is a separate concept (predicted probability
				vs. observed frequency).
			\end{itemize}
		\end{block}
	\end{frame}
	
	\begin{frame}{Q555 --- AUC Comparison}
		\begin{block}{Question}
			Two models have AUC values of 0.91 and 0.73.
			
			What can be concluded?
			\begin{enumerate}
				\item The first model has better discrimination ability
				\item The first model has higher precision
				\item The first model has higher recall
				\item The first model has higher accuracy
			\end{enumerate}
		\end{block}
		
		\begin{exampleblock}{Answer: A}
			The first model has better discrimination ability
		\end{exampleblock}
		
		\begin{block}{Explanation}
			AUC measures \textbf{discrimination / ranking ability}. Since
			\(0.91 > 0.73\), the first model separates positives from negatives better.
			
			\textbf{Why the others are wrong:}
			\begin{itemize}
				\item \textbf{B:} AUC says nothing about precision at a specific threshold.
				\item \textbf{C:} AUC says nothing about recall at a specific threshold.
				\item \textbf{D:} AUC is not accuracy; a model with higher AUC can have
				lower accuracy at a poorly chosen threshold.
			\end{itemize}
			
			\textbf{Key point:} AUC compares \emph{overall ranking quality}, not any single
			threshold metric like precision, recall, or accuracy.
		\end{block}
	\end{frame}
	
	\begin{frame}{Q556 --- Ranking Interpretation}
		\begin{block}{Question}
			An AUC of 0.85 means approximately:
			\begin{enumerate}
				\item 85\% of predictions are correct
				\item 85\% recall
				\item An 85\% chance a randomly chosen positive receives a higher score than a randomly chosen negative
				\item An 85\% false positive rate
			\end{enumerate}
		\end{block}
		
		\begin{exampleblock}{Answer: C}
			An 85\% chance a randomly chosen positive receives a higher score than a randomly chosen negative
		\end{exampleblock}
		
		\begin{block}{Explanation}
			This is the standard \textbf{probabilistic interpretation} of AUC:
			\[
			AUC = P(\text{score of a random positive} > \text{score of a random negative})
			\]
			
			An AUC of 0.85 means that in 85\% of randomly drawn (positive, negative) pairs,
			the model assigns the positive a higher score.
			
			\textbf{Why the others are wrong:}
			\begin{itemize}
				\item \textbf{A:} 85\% accuracy is unrelated to AUC = 0.85.
				\item \textbf{B:} AUC is not recall.
				\item \textbf{D:} AUC is not FPR; FPR is a single-threshold quantity.
			\end{itemize}
			
			\textbf{Memory aid:} AUC = probability of correct ranking of a random pair.
		\end{block}
	\end{frame}
	
	\begin{frame}{Q557 --- Accuracy vs AUC}
		\begin{block}{Question}
			A highly imbalanced fraud dataset contains 99\% non-fraud observations.
			
			Which metric is often more informative than accuracy?
			\begin{enumerate}
				\item AUC
				\item Mean
				\item Variance
				\item R-squared
			\end{enumerate}
		\end{block}
		
		\begin{exampleblock}{Answer: A}
			AUC
		\end{exampleblock}
		
		\begin{block}{Explanation}
			With 99\% non-fraud, a trivial model that predicts ``non-fraud'' for every
			observation achieves \textbf{99\% accuracy} while catching \textbf{zero}
			fraud cases. This is the classic \textbf{class imbalance trap}.
			
			\textbf{AUC is more informative} because:
			\begin{itemize}
				\item It measures discrimination across all thresholds.
				\item It is insensitive to the class distribution.
				\item A random classifier still gives AUC = 0.50, regardless of imbalance.
			\end{itemize}
			
			\textbf{Why the others are wrong:}
			\begin{itemize}
				\item \textbf{B:} Mean is a descriptive statistic, not a classifier metric.
				\item \textbf{C:} Variance is not a classification metric.
				\item \textbf{D:} R-squared is for regression, not classification.
			\end{itemize}
			
			\textbf{Other useful metrics for imbalanced data:} precision, recall, F1,
			and the Precision-Recall curve.
		\end{block}
	\end{frame}
	
	\begin{frame}{Q558 --- ROC Movement}
		\begin{block}{Question}
			Moving along an ROC curve is typically caused by:
			\begin{enumerate}
				\item Changing the classification threshold
				\item Retraining the model
				\item Increasing sample size
				\item Changing the target variable
			\end{enumerate}
		\end{block}
		
		\begin{exampleblock}{Answer: A}
			Changing the classification threshold
		\end{exampleblock}
		
		\begin{block}{Explanation}
			Each point on an ROC curve corresponds to a \textbf{specific classification
				threshold}. As you slide the threshold from 0 to 1:
			\begin{itemize}
				\item Lower threshold \(\rightarrow\) more positives predicted
				\(\rightarrow\) higher TPR and higher FPR (move up-right).
				\item Higher threshold \(\rightarrow\) fewer positives predicted
				\(\rightarrow\) lower TPR and lower FPR (move down-left).
			\end{itemize}
			
			So \textbf{moving along the curve} = changing the threshold on a
			\emph{fixed} model.
			
			\textbf{Why the others are wrong:}
			\begin{itemize}
				\item \textbf{B:} Retraining produces a \emph{different} ROC curve (the
				whole curve shifts), not movement along the same curve.
				\item \textbf{C:} Increasing sample size does not move you along the curve.
				\item \textbf{D:} Changing the target variable changes the problem entirely.
			\end{itemize}
		\end{block}
	\end{frame}
	
	\begin{frame}{Q559 --- GARP Trap}
		\begin{block}{Question}
			Which statement regarding AUC is TRUE?
			\begin{enumerate}
				\item AUC equals accuracy
				\item AUC requires a single threshold
				\item AUC measures discrimination across thresholds
				\item AUC measures precision
			\end{enumerate}
		\end{block}
		
		\begin{exampleblock}{Answer: C}
			AUC measures discrimination across thresholds
		\end{exampleblock}
		
		\begin{block}{Explanation}
			AUC is \textbf{threshold-independent}: it summarizes the model's
			\textbf{discrimination ability} across \emph{all} thresholds by integrating
			the area under the ROC curve.
			
			\textbf{Why the others are wrong:}
			\begin{itemize}
				\item \textbf{A:} AUC \(\neq\) accuracy. They measure different things;
				accuracy is threshold-dependent, AUC is not.
				\item \textbf{B:} AUC requires \emph{no} single threshold---that is its
				main advantage.
				\item \textbf{D:} AUC is not precision. Precision is a single-threshold
				metric.
			\end{itemize}
			
			\textbf{The GARP trap:} The distractors mix up AUC with accuracy, precision,
			and threshold-dependent metrics. Remember: AUC = discrimination across all
			thresholds.
		\end{block}
	\end{frame}
	
\end{document}